\section{Compression over all constituents of the induced unit}
\label{sec:full-I1-compression}

The modular-representation argument in
\cite[Lemma~4.1 and the proof of Theorem~5.1]{SchDimensionBlock}
can retain every simple constituent of $I(\mathbf1)$, rather than only
the Galois orbit of the dimension character.  This gives two further
exclusions in the pivotal and spherical settings.

\begin{proposition}\label{prop:full-I1-bound}
Let $\mathcal C$ be a spherical fusion category, with dimension field
$K=\Q(\dim X:X\in\operatorname{Irr}(\mathcal C))$ and
$G=\operatorname{Gal}(K/\Q)$.  Write $f_E$ for the formal codegrees
of its fusion algebra.  Then $f_E\in K$ and, for every character
$\chi$ of $G$ with conductor $q$,
\[
 \sum_E\left|\sum_{\sigma\in G}
       \chi(\sigma)\sigma(f_E^{-1})\right|^2\leq\frac{|G|}{q}.
\]
\end{proposition}
\begin{proof}
Put $D=\dim\mathcal C$ and $s=S/D$ for the normalized matrix of
$\mathcal Z(\mathcal C)$.  Let $A_E$ be the constituent of
$I(\mathbf1)$ associated to $E$.  It has dimension $D/f_E$ and
twist one, so $f_E\in K$.  If $A_\sigma$ corresponds to
$\sigma(\dim)$, then
$s_{A_\sigma,Y}=\sigma(\dim Y/D)$.
Set $w_\chi=\sum_\sigma\chi(\sigma)e_{A_\sigma}$ and let $P$
project onto the twist-one subspace.  The cited local compression
estimate gives $\|Ps w_\chi\|^2\leq |G|/q$.
Its $A_E$ coordinate is the inner sum in the assertion.  Summing
over these distinct simple coordinates proves the inequality;
their multiplicities in $I(\mathbf1)$ do not repeat coordinates.
\end{proof}

\begin{corollary}\label{cor:r8-quadratic-Casimir}
Let $B=\sum_iN_iN_i^{\mathsf T}$ be the regular Casimir of a fusion
ring.  Suppose the irreducible factors of $\det(xI-B)$ belonging
to a fixed real quadratic field of discriminant $q$ are
$(x^2-t_jx+n_j)^{m_j}$.  A spherical categorification requires
\[
 q\sum_jm_j\frac{t_j^2-4n_j}{n_j^2}\leq1.
\]
\end{corollary}
\begin{proof}
The Casimir eigenvalue attached to $E$ is $\lambda_E=(\dim E)f_E$,
with regular multiplicity $(\dim E)^2$.  The quadratic field lies
in $K$.  If its index in $K$ is $e'$, pull back its nontrivial
character to $G$.  Each conjugate pair of eigenvalues contributes
$2e'^2m_j(t_j^2-4n_j)/n_j^2$ to the preceding bound, whose right
side is $2e'/q$.  Cancelling and using $e'\geq1$ proves the claim.
\end{proof}

\begin{proposition}\label{prop:r8-318-735-scoped}
The rings $R8\_2\_318$ and $R8\_2\_735$ admit no pivotal or
spherical categorification over $\C$.
\end{proposition}
\begin{proof}
For both source tensors, direct multiplication gives
\[
 \det(xI-B)=(x-16)^3(x-8)(x^2-192x+1024)(x^2-16x+32).
\]
Both quadratic factors belong to $\Q(\sqrt2)$, of discriminant
$8$.  Hence the necessary bound would read
\[
 8\left(\frac{192^2-4\cdot1024}{1024^2}
       +\frac{16^2-4\cdot32}{32^2}\right)=\frac54\leq1,
\]
a contradiction.

For completeness, the eight characters of each commutative ring have
formal codegrees
\[16,16,16,8,8+4\sqrt2,8-4\sqrt2,96+64\sqrt2,96-64\sqrt2.\]
Every character in the first six positions vanishes on at least
one basis element.  The remaining two are the Frobenius--Perron
character and its quadratic conjugate, and are duality invariant.
These statements follow by substitution in the source tensor;
the complete exact character tables are supplied in the associated
\texttt{rank8\_full\_I1\_certificate.json}.
Simple pivotal dimensions cannot vanish.  Thus a pivotal structure
would have one of the last two dimension characters and would be
spherical, which has just been excluded.
\end{proof}

These compression arguments alone give scoped exclusions. The independent
module obstruction in Section~\ref{sec:r8-anti-module} now excludes arbitrary
complex categorifications of both rings.

\paragraph{An independent general proof for the four quadratic cases.}
The same estimate gives a second proof of
Proposition~\ref{prop:r8-conductor-four}, without a pseudo-unitary
reduction.  Each of its four rings has Casimir polynomial
\[
 (x-8)^4(x^2-160x+640)(x^2-32x+128).
\]
The sphericalization's dimension field therefore contains both
$\Q(\sqrt{10})$ and $\Q(\sqrt2)$.  Its index $e$ over
$\Q(\sqrt{10})$ is at least two.  The first quadratic factor has
quantity $Q=9/4$ in Corollary~\ref{cor:r8-quadratic-Casimir}.
In the sphericalization's positive-sign sector, formal codegrees are
doubled, so compression requires $eQ/4\leq1$.  But
$eQ/4\geq9/8>1$, a contradiction.
